\section{Ulteriori tecniche di rendering}

Innanzitutto è utile far notare la possibilità di utilizzare approcci ibridi, che
combinano i vantaggi di CSR e SSR per fornire un'esperienza utente ottimale, valutando i compromessi tra interattività e
prestazioni. Un esempio di approccio ibrido è il rendering lato server con \emph{hydration} lato client, una tecnica spiegata nel dettaglio nel prossimo capitolo, \nameref{ch:hydration}.

In aggiunta alle tecniche di rendering discusse precedentemente, esistono altre tecniche che possono essere utilizzate per ottimizzare il rendering di un'applicazione web.
Queste tecniche possono essere utilizzate singolarmente o combinate per ottenere il massimo delle prestazioni e dell'interattività. In questa sezione, esamineremo alcune di queste tecniche\footcite{rendering} e discuteremo i loro vantaggi e svantaggi.

\subsection{Static Site Generation (SSG)}
Il SSG è una tecnica di rendering che consente di generare pagine statiche in
fase di compilazione. Con l'SSG, le pagine vengono generate in anticipo (in fase di build/deploy) e vengono servite ai
visitatori senza la necessità di eseguire codice lato server. Questo approccio consente di ridurre i tempi di
caricamento, di ridurre il carico del server e di migliorare le prestazioni, a discapito della disponibilità di
contenuti aggiornati all'ultima versione.\\
Questa tecnica è particolarmente utile per siti Web con contenuti statici o che non richiedono aggiornamenti frequenti,
ad esempio siti di documentazione o blog.

\subsection{Incremental Static Regeneration (ISR)}
L'ISR è una tecnica di rendering ibrida che combina i vantaggi di SSG
e SSR. Nell'ISR, le pagine vengono generate staticamente, le pagine statiche vengono generate in fase di build, ma
possono essere rigenerate incrementalmente, su richiesta, in modo da riflettere i nuovi contenuti senza dover
ricostruire l'intero sito. Questo approccio permette di mantenere performance elevate, poiché le pagine sono servite
come statiche, ma consente anche di aggiornare contenuti specifici in background a intervalli configurabili.\\
Questa tecnica è particolarmente utile per siti Web con contenuti che cambiano di frequente, ma che non richiedono
aggiornamenti in tempo reale, come e-commerce o siti di news, in cui alcune pagine devono essere rigenerate in base al
traffico e alle necessità, mantenendo la freschezza dei contenuti, ma senza dover rigenerare l'intero sito ad ogni
modifica.

\subsection{Deferred Static Generation (DSG)}
La DSG è una tecnica di rendering che ottimizza ulteriormente il processo
di generazione statica, permettendo di differire la generazione di alcune pagine fino a quando non vengono
effettivamente richieste dagli utenti. 
Quando una pagina \emph{deferred} viene richiesta per la prima volta, viene generata staticamente al volo e servita all'utente. Successivamente, quella pagina viene memorizzata in cache e servita come una normale pagina statica, senza necessità di ulteriore rendering.\\
Questa tecnica è particolarmente utile per siti Web con grandi quantità di contenuti, dove solo una parte del sito è frequentemente visitata, ad esempio forum come Reddit, in cui molti dei contenuti sono vecchi e raramente visitati, ma devono essere disponibili per gli utenti.

